\documentclass{article}
\usepackage{../Self_Style}

\usepackage[normalem]{ulem} %for strike through

\title{UChicago REU Notes Week 5}
\author{Zih-Yu Hsieh}
\date{\today}

\begin{document}
\maketitle
\tableofcontents

The goal for this note is to get a hand of what 2-categories are, both intuitively and some tangible generalization.

\hfil

\section{Motivating Examples}

First, let's consider the most explicit example, the category of small categories, $\Cat$:
\begin{exmp}
    Given two small categories $\mathcal{C},\mathcal{D}\in\Cat$, the homset is the \emph{functor category} $[\mathcal{C},\mathcal{D}]$. Which, the morphism in $[\mathcal{C},\mathcal{D}]$ are the natural transformations between functors. 
\end{exmp}

To keep in mind that we're still doing topology (in most of the cases), let's provide the category of topological spaces $\Top$ as an intuition:

\begin{exmp}
    For each $X,Y\in\Top$, the homset $\Hom_\Top (X,Y)$ is in fact a groupoid: Define morphisms between maps $f,g:X\rightarrow Y$ to be homotopy classes of homotopies (denoted as $[H]:f\Rightarrow g$), with compositions being concatinations of homotopies: If $[H_1]:f\Rightarrow g$ and $[H_2]:g\Rightarrow h$, then $[H_2]\circ_2 [H_1]:=[H_2\cot H_1]$ as follow:
    \[H_2\cdot H_1:X\times I\rightarrow Y,\quad (H_2\cdot H_1)(x,t):= \begin{cases}
        H_1(x,2t) & 0\leq t\leq \frac{1}{2}\\
        H_2(x,2t-1) & \frac{1}{2}<t\leq 1
    \end{cases}\]
    Then, $\Hom_\Top (X,Y)$ forms a groupoid (a category with all morphisms being isomorphisms).
\end{exmp}

\hfil

If associating each small category as topological space, functors as continuous maps, and natural transformations as homotopies, the first example is analogous to the second one. 

An intuitive picture one can have in mind is: In 2-categories, objects are 0-cells, 1-morphisms are 1-cells, and 2-morphisms are 2-cells (serving as transformations between the 1-cells).

\hfil



\newpage

\section{Enriched Categories}

This section follows Emily Riehls' \emph{Categorical Homotopy Theory}.

To properly define notions needed afterward, we need a concrete way of viewing categories enriched over small categories. This more generally follows categories enriched over symmetric monoidal categories, which are categories whose hom functor (or some similarly-behaved functor) are endowed with extra structures relevent to a monoidal category.

\begin{notn}
We'll use $(\mathcal{V},\times,*)$ to denote a symmetric monoidal category here. If monoidally closed, we'll use $\mathcal{V}(\_,\_)$ to denote the internal hom functor. 
\end{notn}

\begin{defn}
    A category $\mathcal{D}$ is \emph{enriched over $\mathcal{V}$}, or called a \emph{$\mathcal{V}$-category}, if it consists of the following data:
    \begin{itemize}
        \item For each pair $x,y\in\mathcal{D}$, there is a \emph{Hom-object} $\mathcal{D}(x,y)\in\mathcal{V}$.
        \item For each $x\in\mathcal{D}$, a \emph{unit morphism} $\id_x:*\rightarrow\mathcal{D}(x,x)$ in $\mathcal{V}$.
        \item For each triple $x,y,z\in\mathcal{D}$, a \emph{composition morphism} $\circ:\mathcal{D}(y,z)\times \mathcal{D}(x,y)\rightarrow \mathcal{D}(x,z)$ in $\mathcal{V}$.
    \end{itemize}
    such that the following diagram commutes when given all $x,y,z,w\in\mathcal{D}$:
    \begin{itemize}
        \item \emph{Associativity of composition}:
        \[\xymatrix{
            \mathcal{D}(z,w)\times\mathcal{D}(y,z)\times\mathcal{D}(x,y) \ar[d]_-{\circ\times 1}\ar[r]^-{1\times \circ} & \mathcal{D}(z,w)\times \mathcal{D}(x,z)\ar[d]^-{\circ}\\
            \mathcal{D}(y,w)\times \mathcal{D}(x,y)\ar[r]_-{\circ} & \mathcal{D}(x,w)
        }\]
        \item \emph{Existence of Two-sided Units:}
        \[\xymatrix{
            \mathcal{D}(x,y)\times \text{*} \ar@{=}[dr] \ar[r]^-{1\times \id_x} & \mathcal{D}(x,y)\times \mathcal{D}(x,x) \ar[d]^-{\circ}\\
            & \mathcal{D}(x,y)
        },\quad\quad \xymatrix{
            \mathcal{D}(y,y)\times \mathcal{D}(x,y)\ar[d]_-{\circ} & \text{*} \times \mathcal{D}(x,y) \ar[l]_-{\id_y\times 1} \ar@{=}[dl]\\
            \mathcal{D}(x,y)
        }\]
    \end{itemize}
\end{defn}

These conditions are mimicing the idea of homsets in a given category. For instance, with this definition by choosing $\mathcal{D}(\_,\_):= \Hom_\mathcal{D}(\_,\_)$, we immediately realize all categories are enriched over $\Sets$. Since more general categories don't have concrete descriptions of objects using sets, we instead need a generalization so that the analogy of ``hom functors'' can land in more general categories. 

%One thing that sorts of questions me, is the existence of $\mathcal{D}$ as only an assignment of objects, not defined directly as a functor. I think in general a functor works just fine. Out of curiosity, does there exist such enrichment example, where $\mathcal{D}$ doesn't form a functor? I think these naturalities should provide a functor structure.  

\hfil

\begin{exmp}
    Given a one-object category $\mu$ enriched over $(\Ab,\tensor_\ZZ,\ZZ)$, the unique object $*$ has a unit homomorphism $\id_*:\ZZ\rightarrow m(*,*)$, a multiplication homomorphism $\cdot:m(*,*)\tensor_\ZZ m(*,*)\rightarrow m(*,*)$. 

    Denote $1_m:=\id_* (1)$, then we precisely have: $x\dot 1_m=x$, $1_m\dot x=x$, and $(x+y)\cdot (z+w)=x\cdot z+x\cdot w+y\cdot z+y\cdot w$, using unit and $\ZZ$-bilinearity of the morphism. So, one-object category enriched over $\Ab$ with tensor product as monoidal product, is a unital ring.
\end{exmp}

\begin{exmp}
    If consider the category of compactly generated spaces $\mathcal{U}$, then the product (i.e. the $k$-ification for original cartesian product) satisfies the following adjunction:
    \[\Hom_\mathcal{U}(X\times K,Y)\cong \Hom_\mathcal{U}(X,Y^K)\]
    where $Y^K:=\Hom_\mathcal{U}(K,Y)$
    Which, it's actually cartesian closed and cocomplete.
\end{exmp}

%I remembered in the book it's mentioned that for general topological space it's not cartesian closed, but do need to think about why.

\begin{exmp}
    If $\mathcal{V}$ is monoidally closed + cocomplete, then its internal hom functor $\mathcal{V}(\_,\_)$ induces a structure where $\mathcal{V}$ is self-enriched.

    The unit morphism is identified by the adjunction $\Hom_\mathcal{V}(*\times v,v)\cong \Hom_\mathcal{V}(*,\mathcal{V}(v,v))$, identified as the image of the natural isomorphism $\gamma_v:*\times v\xrightarrow{\sim}v$, denoted as $\gamma_v^\sharp:*\rightarrow\mathcal{V}(v,v)$. 

    Before building up the composition morphism, we need to consider an \emph{evaluation morphism} $\epsilon:\mathcal{V}(u,v)\times u\rightarrow v$. This is constructed using the same adjunction, together with the pullback of identity $1:\mathcal{V}(u,v)\xrightarrow{\sim}\mathcal{V}(u,v)$, or $\epsilon := 1^\flat$. Then, we define the pushforward of the following morphism:
    \[\phi:\mathcal{V}(v,w)\times \mathcal{V}(u,v)\times u\xrightarrow{\id\times \epsilon}\mathcal{V}(v,w)\times v\xrightarrow{\epsilon}w,\quad \circ := \phi^\sharp:\mathcal{V}(v,w)\times\mathcal{V}(u,v)\rightarrow \mathcal{V}(u,w)\]
    Let's temporarily believe this works out :)
    %(to do: figure out the details for this one)
\end{exmp}

\hfil

\begin{comment}
Let $\mathcal{V}$ be a cartesian closed monoidal category that's also bicomplete (so, categorical product $\times$ is our monoidal product, final object $*\in\mathcal{V}$ serves as unit due to the fact $*\times X\cong X$ by checking universality; it has an internal hom functor $\mathcal{V}(\_,\_):\mathcal{V}^\Op\times \mathcal{V}\rightarrow\mathcal{V}$ generating the adjoint to $\times$ when fixing one variable, while all desired limits / colimits exist).

The goal is to generated something similar to the \emph{category of based objects} analogous to the category of based sets / based spaces over $\Sets$ / $\Top$ respectively, with \emph{smash product} as our monoidal product. Let's first check these two categories for an understanding:
\begin{exmp}
    Given $\Sets$ or $\Top$, the category of based spaces is equivalent to the slice category of the final object $*\in\Top$ (the image of $*$ becomes the basepoint). We have each space $(X,x_0),(Y,y_0)$ satisfies wedge product $X\vee Y:= X\sqcup Y\/x_0\sim y_0$, and smash product $X\wedge Y:= X\times Y\/X\vee Y$ using the identification with basepoint. 
    
    Consider the inclusion $X \rightarrow X\times Y$ by $x\mapsto (x,y_0)$, and $Y\rightarrow X\times Y$ by $y\mapsto (x_0,y)$, this generaets a map $X\sqcup Y \rightarrow X\times Y$. Which, we claim the following fibre coproduct diagram is satisfied:
    \[\xymatrix{
        X\sqcup Y \ar[d] \ar[r] & X\times Y \ar[d] \ar@(r,u)[ddr]\\
        \text{*} \ar@(d,l)[drr] \ar[r] & X\wedge Y \ar@{-->}[dr]\\
        && Z
    }\]
    By sending $*$ to the class of $X\vee Y$ in $X\wedge Y$, since the image of $X\sqcup Y$ in $X\times Y$ is precisely $X\vee Y$, the above diagram commutes. Moreover, if maps into $Z$ makes the above diagram commutes, we must have the image of $X\vee Y\subseteq X\times Y$ be the point specified by $*$, hence must uniquely factor through $X\wedge Y$ using the quotient space / set property.
\end{exmp} 

\begin{con}
    Consider the slice category of the final object $\mathcal{V}_*$, then for any $X,Y\in\mathcal{V}_*$, the product is naturally in the slice category also:
    \[\xymatrix{
        & \text{*} \ar[dl] \ar[dr] \ar@{-->}[d]\\
        X & X\times Y \ar[l] \ar[r] & Y
    }\]
    To mimic the morphism $X\sqcup Y\rightarrow X\times Y$, we consider morphisms $X\rightarrow X\times Y$ (and respectively, for $Y\rightarrow X\times Y$) defined as follow (for $Y$, switch up the $X$ and $Y$):
    \[\xymatrix{
        & X \ar[r] \ar[dl]_-{1_X} \ar@{-->}[d]^-{\exists !} & \text{*} \ar[d]\\
        X & X\times Y \ar[l]\ar[r] & Y
    }\]
    So, after generating these two morphsisms, one naturally induces $X\sqcup Y \rightarrow X\times Y$. So, we can define the \emph{smash product} $X\wedge Y$ as the fibre coproduct as follow:
    \[\xymatrix{
        X\sqcup Y \ar[r] \ar[d] & X\times Y \ar[d]\\
        \text{*} \ar[r] & X\wedge Y
    }\]
    Which, this construction (using the universality) should be associative up to isomorphism.
\end{con}
\end{comment}

\newpage

\section{2-Categories and Bicategories}

Finally, to define a 2-category, there are multiple ways of defining it:
\begin{defn}
    \sout{A category $\mathcal{C}$ is a \emph{2-category}, if it's an $\infty$-category, such that all $3$-morphism is an equivalence, and all $n\geq 3$-morphisms that're parallel are equivalent.}
\end{defn}
If anyone can think in an $\infty$-category manner then they can take this definition, at least I can't...

\hfil

The first (and most intuitive) definition of a 2-category uses enrichment:

\begin{defn}
    A category $\mathcal{C}$ is a \emph{2-category}, if the hom-functor provides an enrichment over $\Cat$, the category of small categories.
\end{defn}
This means, each $X,Y\in\mathcal{C}$ has $\Hom_\mathcal{C}(X,Y)$ actually forms a small category (called a \emph{hom-category}), and the composition of morphisms $\circ:\Hom_\mathcal{C}(Y,Z)\times \Hom_\mathcal{C}(X,Y)\rightarrow\Hom_\mathcal{C}(X,Z)$ is a functor.

Here, we call the elements in $\Hom_\mathcal{C}(X,Y)$ the \emph{1-cells}, and the morphisms within the hom-category are called the \emph{2-cells}.

\hfil

For instance, recalling back to our initial example of $\Top$, and each $\Hom_\Top(X,Y)$ realized as a groupoid using homotopic classes of homotopies as morphisms between maps, this is a well-defined $2$-category structure on $\Top$. Here, I want to emphasize an aspect about the 2-cells, since they have 2 different ``compositions'':

\begin{exmp}
    Fixing two spaces $X,Y$, for any three maps $f,g,h\in \Hom_\Top (X,Y)$, pick two classes of homotopies $[H_1]:f\Rightarrow g$ 
\end{exmp}

\hfil

\hfil



\newpage

\section{Appendix: Categorical Verifications}

\end{document}